\documentclass[11pt,answers]{exam}

\usepackage{etex}
\usepackage{amssymb,amsmath,multicol} %<-- InWorksheetExam1 i also have fancyhdr,

\usepackage[metapost]{mfpic}
\usepackage[pdftex]{graphicx}
\usepackage{tabu}

\usepackage{pst-plot}
\usepackage{pgfplots}
\pgfplotsset{compat=1.9}

\usepackage{tikz}
\usepackage{tkz-2d}
\usepackage{tkz-base}
\usetikzlibrary{calc}
\usetikzlibrary{arrows}

\usepackage{systeme}

\usepackage[inline]{enumitem}
\usepackage{refcount}%<-- non in WorksheetExam1

\usepackage{pstricks-add,pst-eucl}
\usepackage{systeme}
\usepackage{setspace}
\usepackage{multicol}


\usepackage[inline]{enumitem}   
\makeatletter
% This command ignores the optional argument for itemize and enumerate lists
\newcommand{\inlineitem}[1][]{%
\ifnum\enit@type=\tw@
    {\descriptionlabel{#1}}
  \hspace{\labelsep}%
\else
  \ifnum\enit@type=\z@
       \refstepcounter{\@listctr}\fi
    \quad\@itemlabel\hspace{\labelsep}%
\fi}
\makeatother


\def\f{x+1} \def\g{-x/3+2}  \def\h{-x+3}

\newcommand{\vasymptote}[2][]{
    \draw [densely dashed,#1] ({rel axis cs:0,0} -| {axis cs:#2,0}) -- ({rel axis cs:0,1} -| {axis cs:#2,0});
}

\addpoints
%\printanswers
\noprintanswers

\opengraphsfile{Q5c_Fa18}

\begin{document}
\extrawidth{-0.3in}
\pagestyle{headandfoot}

\setlength{\hoffset}{-.25in}

\extraheadheight{-.3in}
\runningheadrule
\firstpageheader{\bfseries {Precalculus}}{ \bfseries {Quiz 5 }}{\bfseries {10/16/18}} 

\begin{center}
	This quiz has \numquestions\ questions, for a total of \numpoints\
	points and \numbonuspoints\ bonus points.
\end{center}


\firstpagefooter{} %%&&CHANGED
                {}
                {%Points earned: \hbox to 0.5in{\hrulefill}
                % out of  \pointsonpage{\thepage} points
                }
                 
						

\vspace*{0.1cm}
\hbox to \textwidth { \scshape {Name:} \enspace\hrulefill}
\vspace{0.1cm}




\pointpoints{point}{points}

\begin{questions}


\addpoints

\question The quadratic function $f(x)$ is given by: $\displaystyle f(x) = -(x+2)^2+2$. 

\begin{parts}
	\part[2] Write the range of  $\displaystyle f(x)$ in interval form. Show evidence of Your thinking.
	
	\fillwithdottedlines{2cm}
	
	\part[1] As $x\to\infty$, $f(x)\to$ \dotfill
	\part[1] As $x\to-\infty$, $f(x)\to$ \dotfill
	
	
\end{parts}

\question[1] True or False? The polynomial function $P$has degree 4, one local maximum, and no local minima.
\begin{oneparchoices}
	\choice Possibly true \choice Impossible
\end{oneparchoices}

\question The polynomial function $Q(x)$ is given by: 
$\displaystyle Q(x) =-2 (x-1)^2(x+1)^3$.

\begin{parts}
	\part[1] As $x\to\infty$, $Q(x)\to$ \dotfill
	\part[1] As $x\to -\infty$, $Q(x)\to$ \dotfill
	%\part[1] The $y$-intercept of $Q(x)$ is: \dotfill
	\part[2] Find all the $x$-intercepts and state their multiplicity.
	
\fillwithdottedlines{3cm}

\end{parts}

\bonusquestion[2]  Draw a possible graph of a polynomial $P(x)$ of even degree, which satisfies all of the following:
\begin{itemize}
	\item There are exactly four $x$-intercepts:  $(1,0)$, $(2,0)$, $(3,0)$ and $(4,0)$;
	\item The graph bounces off the $x$-axis at $(2,0)$;
	\item The graph bounces off the $x$-axis at $(3,0)$.
	
\end{itemize}   

Make sure to draw and label the $x$ and $y$-axes and include units.

\begin{tikzpicture}[thick,scale=0.9, every node/.style={scale=0.6}]

\draw[gray!70, thin, step=0.5] (-1,-3) grid (\linewidth-\pgflinewidth,3);

\end{tikzpicture} 

\end{questions}

\end{document}                 